\makesubsectionpage{Support Vector Machines}
\makesubsubsectionpage{Grundidee}

\begin{frame}{Support Vector Machines: Grundidee}
\begin{columns}
\column{0.5\textwidth}
Klassifikatoren unterteilen den Eingaberaum in \textbf{Bereiche}, die
unterschiedlich klassifiziert werden.

\vspace{0.4em}
\begin{itemize}
  \item Die Grenzen zwischen den Bereichen sind die
        \textbf{Entscheidungsgrenzen}.
  \item Diese haben wir bei den bisherigen Klassifikatoren bereits
        betrachtet.
\end{itemize}

\vspace{0.4em}
{\footnotesize\itshape Beispiel rechts: Entscheidungsgrenze eines
Entscheidungsbaumes -- achsenparallele Stufen.}

\column{0.45\textwidth}
\centering
\includesvg[width=\linewidth]{figures/scatter_baum_beschraenkt.svg}
\end{columns}
\source{\cite{seemann2024mlaz}}
\end{frame}





\begin{frame}{Support Vector Machines: Large Margin}
\begin{columns}
\column{0.50\textwidth}
Die Klassifikation wird \textbf{robuster}, wenn der Abstand zwischen
Entscheidungsgrenze und Datenpunkten \textbf{groß} ist.

\vspace{0.4em}
\begin{itemize}
  \item SVMs sind \textbf{Large-Margin-Klassifikatoren}.
  \item Sie platzieren eine Trennebene mit \textbf{maximalem Abstand} zu den
        Datenpunkten.
  \item Diese Trennebene ist die \textbf{Entscheidungsgrenze}.
\end{itemize}

\textbf{Der zentrale Unterschied zu anderen Klassifikatoren:} Die Grenze
entsteht \textbf{nicht zufällig} als Nebenprodukt des Trainings. SVMs legen
sie \textbf{ganz gezielt} so, dass der \textbf{Abstand} zu den nächsten
Datenpunkten \textbf{maximal} wird.

\column{0.48\textwidth}
\centering
\resizebox{\linewidth}{!}{%
\begin{tikzpicture}[
  gut/.style={circle, fill=SecondaryColor, draw=SecondaryTitle,
              very thick, minimum size=6pt, inner sep=0pt},
  schlecht/.style={circle, fill=AccentColor, draw=AccentTitle,
                   very thick, minimum size=6pt, inner sep=0pt}
]
% Rahmen
\draw[thick] (0,0) rectangle (7,7);

% Trennebene (fallende Diagonale)
\draw[AlertTitle, line width=1.5pt] (0,6.6) -- (7,0.4);
\node[AlertTitle, font=\footnotesize\itshape, anchor=south west]
  at (-0.2,7.1) {Trennebene};

% gestrichelte Margin-Linien (parallel)
\draw[SecondaryTitle, dashed] (0.5,7) -- (7,1.2);
\draw[SecondaryTitle, dashed] (0,5.8) -- (6.6,0.0);

% blaue Punkte (oben rechts)
\foreach \x/\y in {4.2/6.2, 3.7/5.2, 4.7/5.0, 4.2/4.4, 5.0/4.0,
                   5.0/3.0, 5.7/4.6, 6.0/3.4, 4.6/3.5} {
  \node[gut] at (\x,\y) {};
}
% orange Punkte (unten links)
\foreach \x/\y in {2.6/3.5, 1.8/2.6, 3.0/2.6, 2.4/1.9, 1.4/1.2,
                   3.0/1.2, 2.2/0.6} {
  \node[schlecht] at (\x,\y) {};
}

% Abstandspfeile (Margin)
\draw[<->, thick] (2.6,3.5) -- (3.0,3.9);
\draw[<->, thick] (5.0,3.0) -- (4.6,2.6);

\end{tikzpicture}}
\end{columns}
\source{\cite{seemann2024mlaz}}
\end{frame}


\begin{frame}{Mathematische Grundlage}
    \footnotesize
\begin{columns}
\column{0.55\textwidth}

Im dreidimensionalen Raum ist eine \textbf{Ebene in Koordinatenform}
definiert durch:
\[
  a x_1 + b x_2 + c x_3 = d
\]
Alle Punkte $(x_1, x_2, x_3)$, die diese Gleichung erfüllen, liegen in der
Ebene. Dabei sind:

\vspace{0.3em}
\begin{itemize}
  \item $a, b, c$: die Komponenten des \textbf{Normalenvektors}
        $\mathbf{n} = (a, b, c)^T$. Er steht \textbf{senkrecht} auf der
        Ebene und bestimmt ihre Ausrichtung.
  \item $d$: bestimmt die \textbf{Lage} der Ebene -- der Abstand vom
        Ursprung (entlang $\mathbf{n}$), multipliziert mit der Länge von
        $\mathbf{n}$.
\end{itemize}

\column{0.4\textwidth}
\centering
\includesvg[width=\linewidth]{figures/svm_normalenvektor.svg}\\[4pt]
{\footnotesize\itshape Die Ebene $2x_1 - 3x_2 + x_3 = 6$ mit Normalenvektor
$(2, -3, 1)^T$.}
\end{columns}
\source{\cite{Frochte2020}}
\end{frame}


\begin{frame}{Support Vector Machines: Mathematische Grundlage}
\begin{columns}
\column{0.4\textwidth}
\centering
\includesvg[width=\linewidth]{figures/svm_normalenvektor.svg}

\column{0.55\textwidth}
\footnotesize
Die Ebene trennt Punkte im Raum in \textbf{„vor"} und \textbf{„hinter"} der
Ebene. Jeder Punkt $P$ liegt in einer Parallelebene zur unseren.

\vspace{0.3em}
Über den Abstand der Parallelebene durch $P$:
\begin{itemize}
  \item Abstand $> d$: $P$ liegt \textbf{vor} der Ebene.
  \item Abstand $< d$: $P$ liegt \textbf{dahinter}.
\end{itemize}

\vspace{0.3em}
\textbf{Berechnung der Lage von $P$:}
\[
  g(x_1, x_2, x_3) = a x_1 + b x_2 + c x_3 - d
\]
\begin{itemize}
  \item $g(x_1, x_2, x_3) > 0$: Punkt \textbf{vor} der Ebene.
  \item $g(x_1, x_2, x_3) < 0$: Punkt \textbf{hinter} der Ebene.
\end{itemize}
\end{columns}
\source{\cite{Frochte2020}}
\end{frame}

\begin{frame}{Support Vector Machines: Hyperebene}
\begin{columns}
\column{0.60\textwidth}
\begin{itemize}
  \item Beim maschinellen Lernen arbeiten wir oft in \textbf{hochdimensionalen}
        Datenräumen: der Iris-Datensatz hat 4 Dimensionen, der
        Heart-Disease-Datensatz nach Aufbereitung 20.
  \item Die Ebenengleichung ist nicht an den 3D-Raum gebunden. Eine
        \textbf{Hyperebene} verallgemeinert die Ebene auf $n$ Dimensionen:
        \[
          w_1 x_1 + w_2 x_2 + \dots + w_n x_n = d
        \]
  \item Der Normalenvektor $(w_1, w_2, \dots, w_n)^T$ bestimmt die
        \textbf{Orientierung}.
\end{itemize}

\column{0.35\textwidth}
\begin{blockC1}{Anschauung}
\begin{itemize}
  \item Im \textbf{3D}-Raum: eine gewöhnliche \textbf{Ebene}.
  \item Im \textbf{2D}-Raum: eine \textbf{Gerade}.
\end{itemize}
\end{blockC1}
\end{columns}
\source{\cite{Frochte2020}}
\end{frame}


\begin{frame}{SVM: Optimale Trennung -- Umformung}
\begin{columns}
\column{0.50\textwidth}
\footnotesize
Wir formen die Hyperebenengleichung um, um die optimale Lage zu finden:

\vspace{0.3em}
\begin{itemize}
  \item Setze $w_0 = -d$ und bringe es auf die andere Seite. Für einen Punkt
        $\mathbf{x} = (x_1, \dots, x_n) \in \mathbb{R}^n$:
        \[
          g(\mathbf{x}) = w_0 + w_1 x_1 + \dots + w_n x_n
        \]
  \item Fasse die $w_i$ im Normalenvektor
        $\mathbf{w} = (w_1, \dots, w_n) \in \mathbb{R}^n$ zusammen:
        \[
          \boxed{\,g(\mathbf{x}) = w_0 + \mathbf{w} \cdot  \mathbf{x}\,}
        \]
\end{itemize}

\column{0.45\textwidth}
\centering
\resizebox{\linewidth}{!}{%
\begin{tikzpicture}[>=Stealth,
  gut/.style={circle, fill=SecondaryColor, draw=SecondaryTitle,
              very thick, minimum size=5pt, inner sep=0pt},
  schlecht/.style={circle, fill=AccentColor, draw=AccentTitle,
                   very thick, minimum size=5pt, inner sep=0pt}
]
% Achsen
\draw[->, thick] (-0.3,0) -- (6,0) node[right] {$x_1$};
\draw[->, thick] (0,-0.3) -- (0,5.5) node[above] {$x_2$};
% Gerade g
\draw[line width=1.2pt] (-0.3,2.0) -- (5.8,4.6) node[right, font=\bfseries] {$g$};
% Achsenabschnitt w0
\draw[very thick, dashed, ->, draw=PrimaryTitle] (0,0) -- (0,2.13);
\node[anchor=east, font=\footnotesize] at (0,2.3) {$(0,w_0)$};
% orange Punkte (Klasse 1, oben links)
\node[schlecht] at (2.85,4.85) {};
\node[schlecht] at (1.7,4.7) {};
\node[schlecht] at (0.5,3.4) {};
% blaue Punkte (Klasse 2, unten rechts)
\node[gut] at (4.3,2.9) {};
\node[gut] at (4.9,2.4) {};
\node[gut] at (3.6,1.5) {};
\node[gut] at (2.2,1.0) {};

% Normalenvektor w
\coordinate (fuss) at (3.4,3.55);
\draw[thick, dashed, ->,draw=PrimaryTitle] (fuss) -- ++(-0.55,1.3);
\node[font=\footnotesize] at (3.5,4.3) {$\mathbf{w}$ };
%%\node[font=\footnotesize\bfseries] at (2.6,4.9) {$\mathbf{v}$};
\end{tikzpicture}}
\end{columns}
\source{\cite{Frochte2020}}
\end{frame}




\begin{frame}{SVM: Herleitung der Abstandsformel (1)}
\begin{columns}

\column{0.40\textwidth}
\centering
\resizebox{\linewidth}{!}{%
\begin{tikzpicture}[>=Stealth,
  gut/.style={circle, fill=SecondaryColor, draw=SecondaryTitle,
              very thick, minimum size=5pt, inner sep=0pt},
  schlecht/.style={circle, fill=AccentColor, draw=AccentTitle,
                   very thick, minimum size=5pt, inner sep=0pt}
]
\draw[->, thick] (-0.3,0) -- (6,0) node[right] {$x_1$};
\draw[->, thick] (0,-0.3) -- (0,5.5) node[above] {$x_2$};
\draw[line width=1.2pt] (-0.3,2.0) -- (5.8,4.6) node[right, font=\bfseries] {$g$};
% Fusspunkt x0 auf der Ebene
\coordinate (x0) at (3.4,3.55);
\node[circle, fill=black, minimum size=3pt, inner sep=0pt] at (x0) {};
\node[font=\footnotesize, anchor=north] at (3.5,3.45) {$\mathbf{x}_0$};
% Punkt x
\coordinate (x) at (2.85,4.85);
\node[schlecht] at (x) {};
\node[font=\footnotesize, anchor=south] at (2.85,5.0) {$\mathbf{x}$};
% Abstand r entlang der Normalen
\draw[thick, dashed, ->, draw=PrimaryTitle] (x0) -- (x);
\node[font=\footnotesize, draw=PrimaryTitle] at (2.75,4.15) {$r$};
% weitere Punkte
\node[schlecht] at (1.5,4.5) {};
\node[schlecht] at (0.5,3.4) {};
\node[gut] at (4.3,2.9) {};
\node[gut] at (4.9,2.4) {};
\node[gut] at (3.6,1.5) {};
\end{tikzpicture}}

\column{0.60\textwidth}
\par\bigskip
\footnotesize
\textbf{Wie weit ist ein Punkt $\mathbf{x}$ von der Ebene entfernt?}
\begin{itemize}
  \item Wir nehmen an, wir kennen Fußpunkt $\mathbf{x}_0$ \textbf{auf} der Hyperebene. Da $\mathbf{x}_0$ auf $g$ liegt, 
        gilt\ $g(\mathbf{x}_0) = 0$.
  \item Der kürzeste Weg von $\mathbf{x}$ zur Ebene verläuft entlang der
        \textbf{Normalen} $\mathbf{w}$, denn sie steht senkrecht auf der
        Ebene.
  \item $\mathbf{x}$ erreicht man also von $\mathbf{x}_0$ aus über ein Stück in
        Normalenrichtung:
        \begin{equation}
          \mathbf{x} = \mathbf{x}_0
                     + r\cdot\frac{\mathbf{w}}{\lVert\mathbf{w}\rVert}
          \label{eq:svm-abstand}
        \end{equation}
        Dabei ist $r$ der gesuchte Abstand und
        $\frac{\mathbf{w}}{\lVert\mathbf{w}\rVert}$ der \textbf{normierte}
        Normalenvektor (Länge 1).
    \item \textbf{Einsetzen} von \eqref{eq:svm-abstand} in $g$ liefert:
        \[
          g(\mathbf{x}) = \mathbf{w} \cdot \!\left(\mathbf{x}_0
                        + r\frac{\mathbf{w}}{\lVert\mathbf{w}\rVert}\right) + w_0
        \]
\end{itemize}

\end{columns}
\source{\cite{Frochte2020}}
\end{frame}


\begin{frame}{SVM: Herleitung der Abstandsformel (2)}
\begin{columns}

\column{0.47\textwidth}
\footnotesize
\begin{itemize}
  \item Ausmultiplizieren:
        \[
          g(\mathbf{x}) = \underbrace{\mathbf{w} \cdot \mathbf{x}_0 + w_0}_{= \, 0}
                        + \; r\cdot\frac{\mathbf{w} \cdot \mathbf{w}}{\lVert\mathbf{w}\rVert}
        \]
  \item Erster Term $= 0$ (Punkt auf Ebene), $\mathbf{w} \cdot \mathbf{w} =
        \lVert\mathbf{w}\rVert^2$:
        \[
          g(\mathbf{x}) = r\cdot\lVert\mathbf{w}\rVert
        \]
  \item Nach $r$ auflösen:
        \[
          \boxed{\,r = \frac{g(\mathbf{x})}{\lVert\mathbf{w}\rVert}\,}
        \]
\end{itemize}

\column{0.48\textwidth}
\footnotesize
\begin{itemize}
  \item $g(\mathbf{x})$ wächst \textbf{proportional} zur Länge von
        $\mathbf{w}$.
  \item Division durch $\lVert\mathbf{w}\rVert$ liefert die \textbf{echte}
        geometrische Distanz.
  \item \textbf{Klassifikation}: nur Vorzeichen nötig -- Normierung
        entfällt.
  \item \textbf{Abstand}: Normierung wird gebraucht.
\end{itemize}

\end{columns}
\source{\cite{Frochte2020}}
\end{frame}


\begin{frame}{SVM: Optimale Trennung -- Klassifikation}
\begin{columns}

\column{0.47\textwidth}
\centering
\resizebox{\linewidth}{!}{%
\begin{tikzpicture}[>=Stealth,
  gut/.style={circle, fill=SecondaryColor, draw=SecondaryTitle,
              very thick, minimum size=5pt, inner sep=0pt},
  schlecht/.style={circle, fill=AccentColor, draw=AccentTitle,
                   very thick, minimum size=5pt, inner sep=0pt}
]
\draw[->, thick] (-0.3,0) -- (6,0) node[right] {$x_1$};
\draw[->, thick] (0,-0.3) -- (0,5.5) node[above] {$x_2$};
\draw[line width=1.2pt] (-0.3,2.0) -- (5.8,4.6) node[right, font=\bfseries] {$g$};
% Klasse A (orange, oben links)
\node[schlecht] at (2.85,4.85) {};
\node[schlecht] at (1.7,4.7) {};
\node[schlecht] at (0.5,3.4) {};
\node[schlecht, font=\footnotesize] at (1.3,4.0) {};
\node[font=\footnotesize, anchor=south] at (1.3,4.15) {A};
% Klasse B (blau, unten rechts)
\node[gut] at (4.3,2.9) {};
\node[gut] at (4.9,2.4) {};
\node[gut] at (3.6,1.5) {};
\node[gut] at (2.2,1.0) {};
\node[font=\footnotesize, anchor=north] at (4.6,2.6) {B};
\end{tikzpicture}}

\column{0.48\textwidth}
\footnotesize
Für die \textbf{Klassifikation} zählt nur das \textbf{Vorzeichen} von
$g(\mathbf{x})$ -- es zeigt die Seite der Trennebene an. Die Normierung mit
$\lVert \mathbf{w} \rVert$ kann daher \textbf{entfallen}:
\[
  \begin{aligned}
    \mathbf{w} \cdot \mathbf{x} + w_0 &> 0 \;\Rightarrow\; \text{Klasse A}\\
    \mathbf{w} \cdot \mathbf{x} + w_0 &< 0 \;\Rightarrow\; \text{Klasse B}
  \end{aligned}
\]

Ein Punkt \textbf{auf} der Ebene ($g(\mathbf{x}) = 0$) liegt genau auf der
Entscheidungsgrenze.

\end{columns}
\source{\cite{Frochte2020}}
\end{frame}


\begin{frame}{SVM: Optimale Trennung}
\begin{columns}

\column{0.5\textwidth}
\footnotesize
Für einen \textbf{größeren Abstand} der Punkte zur Trennebene fordert man Mindestabstand $1$:
\[
  \begin{aligned}
    \mathbf{w} \cdot\mathbf{x} + w_0 &\geq +1 \;\Rightarrow\; \text{Klasse A}\\
    \mathbf{w} \cdot\mathbf{x} + w_0 &\leq -1 \;\Rightarrow\; \text{Klasse B}
  \end{aligned}
\]

Mit $y = +1$ (Klasse A) und $y = -1$ (Klasse B) schreibt man kompakter:
\[
  y\cdot(\mathbf{w} \cdot\mathbf{x} + w_0) \geq +1
\]

Für die \textbf{Optimierung} maximiert man den Abstand $\rho$ mit:
\[
  \frac{y\cdot(\mathbf{w} \cdot\mathbf{x} + w_0)}{\lVert\mathbf{w}\rVert} \geq \rho
\]
für alle Trainingspunkte $\mathbf{x}$.

\column{0.45\textwidth}
\centering
\resizebox{\linewidth}{!}{%
\begin{tikzpicture}[
  gut/.style={circle, fill=SecondaryColor, draw=SecondaryTitle,
              very thick, minimum size=6pt, inner sep=0pt},
  schlecht/.style={circle, fill=AccentColor, draw=AccentTitle,
                   very thick, minimum size=6pt, inner sep=0pt}
]
% Rahmen
\draw[thick] (0,0) rectangle (7,7);

% Trennebene (fallende Diagonale)
\draw[AlertTitle, line width=1.5pt] (0,6.6) -- (7,0.4);

% gestrichelte Margin-Linien (parallel)
\draw[SecondaryTitle, dashed] (0.5,7) -- (7,1.2);
\draw[SecondaryTitle, dashed] (0,5.8) -- (6.6,0.0);

% blaue Punkte (oben rechts)
\foreach \x/\y in {4.2/6.2, 3.7/5.2, 4.7/5.0, 4.2/4.4, 5.0/4.0,
                   5.0/3.0, 5.7/4.6, 6.0/3.4, 4.6/3.5} {
  \node[gut] at (\x,\y) {};
}
% orange Punkte (unten links)
\foreach \x/\y in {2.6/3.5, 1.8/2.6, 3.0/2.6, 2.4/1.9, 1.4/1.2,
                   3.0/1.2, 2.2/0.6} {
  \node[schlecht] at (\x,\y) {};
}

% Abstandspfeile (Margin) mit rho beschriftet
\draw[<->, thick] (2.6,3.5) -- (3.0,3.9);
\node[font=\large] at (2.55,3.95) {$\rho$};
\draw[<->, thick] (5.0,3.0) -- (4.6,2.6);
\node[font=\large] at (5.05,2.55) {$\rho$};

\end{tikzpicture}}

\end{columns}
\source{\cite{Frochte2020}}
\end{frame}

\begin{frame}{SVM: Eindeutigkeit und Margin}
\begin{columns}

\column{0.47\textwidth}
\footnotesize
\begin{itemize}
  \item Der Normalenvektor $\mathbf{w}$ ist \textbf{nicht eindeutig}:
        Vielfache von $\mathbf{w}$ beschreiben dieselbe Trennebene.
  \item Damit ist auch $\rho$ nicht eindeutig, da es von $\mathbf{w}$
        abhängt.
  \item Wir \textbf{fixieren} die Skalierung über
        $\rho \cdot \lVert \mathbf{w} \rVert = 1$, also:
        \[
          \rho = \frac{1}{\lVert \mathbf{w} \rVert}
        \]
  \item Eingesetzt ergibt sich für alle Trainingspunkte $\mathbf{x}$:
        \[
          \frac{y\cdot(\mathbf{w}\cdot\mathbf{x} + w_0)}{\lVert \mathbf{w} \rVert}
          \geq \frac{1}{\lVert \mathbf{w} \rVert}
        \]
\end{itemize}

\column{0.48\textwidth}
\footnotesize
\begin{blockC1}{\footnotesize   Margin}
\begin{itemize}
  \item Die rechte Seite ist der \textbf{Margin} -- der Abstand der
        Trennebene zu den nächstgelegenen Trainingspunkten.
  \item \textbf{Kernidee:} Margin maximieren bedeutet
        $\lVert \mathbf{w} \rVert$ \textbf{minimieren}.
\end{itemize}
\end{blockC1}

\begin{blockC2}{\footnotesize Support Vectors}
\begin{itemize}
  \item Nur wenige Punkte bestimmen die Lage der Ebene: die ihr am nächsten
        liegenden.
  \item Diese heißen \textbf{Support Vectors}.
\end{itemize}
\end{blockC2}

\end{columns}
\source{\cite{Frochte2020}}
\note{
Die Fixierung $\rho \cdot \lVert\mathbf{w}\rVert = 1$ ist keine geometrische
Umrechnung, sondern eine Eichung: Da jedes Vielfache von $\mathbf{w}$
dieselbe Ebene beschreibt, waehlt man gezielt das $\mathbf{w}$, bei dem der
naechstgelegene Punkt genau den Funktionswert $g = 1$ erhaelt. Ohne diese
Festlegung waere die Maximierung von $\rho = g / \lVert\mathbf{w}\rVert$
nicht wohldefiniert, da man Zaehler und Nenner beliebig mitskalieren
koennte. Mit fixiertem Zaehler wird der Margin zu $\rho = 1/\lVert\mathbf{w}
\rVert$; ihn zu maximieren ist damit gleichbedeutend mit der Minimierung von
$\lVert\mathbf{w}\rVert$ bzw. $\tfrac{1}{2}\lVert\mathbf{w}\rVert^2$. Das ist
der Schritt, der die SVM-Optimierung handhabbar macht.
}
\end{frame}


%%%%%%%%%%%%%%%%%%%%%%%%%%%%%%%%%%%%%%%%%%%%%%%%%%%%%%%%%
\begin{frame}{SVM: Das quadratische Optimierungsproblem}
\begin{columns}

\column{0.60\textwidth}
\begin{itemize}
  \item Die Maximierung von $\frac{1}{\lVert\mathbf{w}\rVert}$ ist für
        Optimierer \textbf{ungünstig} (Wurzel im Nenner).
  \item Mit \textbf{quadratischen Problemen} kommen sie dagegen gut zurecht.
  \item $\frac{1}{\lVert\mathbf{w}\rVert}$ maximieren $\Leftrightarrow$
        $\lVert\mathbf{w}\rVert^2$ \textbf{minimieren}.
  \item Die Trennung der Klassen wird als \textbf{Nebenbedingung}
        formuliert:
        \[
          \min \tfrac{1}{2}\lVert\mathbf{w}\rVert^2 \quad\text{u.~d.~N.}\quad
          y \cdot (w_0 + \mathbf{w}\cdot\mathbf{x}) \geq +1
        \]
\end{itemize}

\column{0.35\textwidth}
\begin{blockC1}{Quadratisches Problem}
Die Zielfunktion ist \textbf{quadratisch} (Variablen in zweiter Potenz),
die Nebenbedingungen \textbf{linear}. Solche Probleme sind \textbf{konvex}:
Es gibt genau ein globales Minimum.
\end{blockC1}
\end{columns}
\source{\cite{scholkopf2002learning}}

\note{u.n.N.: Unter der Nebenbedingung}

\end{frame}

\begin{frame}{SVM: Nebenbedingungen und Faktor $\tfrac{1}{2}$}
\begin{columns}

\column{0.60\textwidth}

\textbf{Was sind Nebenbedingungen?}
\begin{itemize}
  \item Bei einer Optimierung minimiert man eine \textbf{Zielfunktion}
        (hier $\tfrac{1}{2}\lVert\mathbf{w}\rVert^2$).
  \item \textbf{Nebenbedingungen} schränken die zulässigen Lösungen ein --
        nur Lösungen, die sie erfüllen, sind erlaubt.
  \item Hier: $y \cdot (w_0 + \mathbf{w}\cdot\mathbf{x}) \geq +1$ für
        \textbf{jeden} Trainingspunkt.
  \item Das erzwingt, dass die Ebene beide Klassen korrekt trennt.
\end{itemize}

\column{0.35\textwidth}
\begin{blockC1}{Warum $\tfrac{1}{2}$?}
Der Faktor $\tfrac{1}{2}$ ist reine Bequemlichkeit. Wir müssen das nachher ableiten. Dann fällt der Faktor $2$ aus $\lVert\mathbf{w}\rVert^2$ weg.
\end{blockC1}
\end{columns}
\source{\cite{scholkopf2002learning}}
\end{frame}




%%%%%%%%%%%%%%%%%%%%%%%%%%%%%%%%%%%%%%%%%%%%%%%%%%%%%%
\begin{frame}{SVM: Lagrange-Multiplikatoren}
\begin{columns}

\column{0.60\textwidth}
\footnotesize
\begin{itemize}
  \item Optimierung unter \textbf{Ungleichungs}-Nebenbedingungen erfordert
        \textbf{Lagrange-Multiplikatoren} und die
        \textbf{Karush-Kuhn-Tucker}-Bedingungen.
  \item Man koppelt Zielfunktion und Nebenbedingungen über neue Variablen
        $\lambda_i$ in einer Funktion:
        \[
          L(\mathbf{w}, w_0, \boldsymbol{\lambda}) =
          \tfrac{1}{2}\mathbf{w}\cdot\mathbf{w}
          + \sum_{i=1}^{n} \lambda_i\bigl(1 - y_i \cdot (w_0 + \mathbf{w}\cdot\mathbf{x}_{i})\bigr)
        \]
  \item $L$ hängt von den zu optimierenden Größen \textbf{$\mathbf{w}$},
        \textbf{$w_0$} und den Multiplikatoren
        \textbf{$\boldsymbol{\lambda} = (\lambda_1, \dots, \lambda_n)$} ab.
  \item Dabei ist $\mathbf{x}_{i}$ der Merkmalsvektor des $i$-ten
        Trainingsdatensatzes, $y_i \in \{+1, -1\}$ sein Klassenlabel.
\end{itemize}

\column{0.35\textwidth}
\begin{blockC1}{\footnotesize Idee}
\footnotesize
Jede Nebenbedingung wird nach Null umgestellt und mit einem eigenen
$\lambda_i$ gewichtet an die Zielfunktion angehängt.
\end{blockC1}
\end{columns}
\source{\cite{scholkopf2002learning}}
\end{frame}

%%%%%%%%%%%%%%%%%%%%%%%%%%%%%%%%%%%%%%%%%%%%%%%%%%%%%%

\begin{frame}{SVM: Karush-Kuhn-Tucker-Bedingungen}
\begin{columns}

\column{0.35\textwidth}
\begin{blockC2}{Notwendige Bedingungen}

Die KKT-Bedingungen sind \textbf{notwendige} Bedingungen: Jede optimale
Lösung erfüllt sie. Sie liefern die Gleichungen, aus denen sich die Lösung
bestimmen lässt.
\end{blockC2}

\column{0.60\textwidth}

\begin{itemize}
  \item Lagrange-Multiplikatoren sind für \textbf{Gleichungs}-Neben\-%
        bedingungen gedacht.
  \item Bei der SVM sind die Nebenbedingungen aber \textbf{Ungleichungen}
        ($y \cdot (w_0 + \mathbf{w}\cdot\mathbf{x}) \geq 1$).
  \item Dann liefert nicht mehr das Nullsetzen der Ableitung die Lösung.
  \item Stattdessen muss jede Lösung die \textbf{KKT-Bedingungen} erfüllen
        -- eine Verallgemeinerung des Lagrange-Ansatzes auf Ungleichungen.
\end{itemize}


\end{columns}
\source{\cite{scholkopf2002learning}}
\end{frame}

\begin{frame}{SVM: Die KKT-Bedingungen}
\small

\begin{columns}

\column{0.47\textwidth}
Für unser \textbf{spezifisches Problem} lauten die KKT-Bedingungen ($*$
markiert das Optimum):
\[
  \begin{aligned}
    \tfrac{\partial L}{\partial \mathbf{w}} = 0,\quad
    \tfrac{\partial L}{\partial w_0} &= 0\\[4pt]
    \lambda_i^{*}\bigl(1 - y_i \cdot (\mathbf{w}^{*}\!\cdot\mathbf{x}_i + w_0^{*})\bigr) &= 0\\[2pt]
    1 - y_i \cdot (\mathbf{w}^{*}\!\cdot\mathbf{x}_i + w_0^{*}) &\leq 0\\[2pt]
    \lambda_i^{*} &\geq 0
  \end{aligned}
\]

\column{0.48\textwidth}
\begin{blockC1}{\small Deutung}
\begin{itemize}
  \item Die erste Zeile ist die \textbf{Stationarität}: Im Optimum werden
        die Ableitungen null.
  \item Die dritte Zeile ist die umgestellte \textbf{Nebenbedingung}.
  \item Die zweite Zeile ist ein Produkt: null, wenn $\lambda_i^{*} = 0$
        \textbf{oder} der Punkt auf der Grenze ($\pm 1$) liegt.
  \item Letzteres gilt nur für die \textbf{Support Vectors}.
  \item $\Rightarrow$ Sehr viele $\lambda_i^{*}$ sind \textbf{null}.
\end{itemize}
\end{blockC1}
\end{columns}
\source{\cite{scholkopf2002learning}}
\end{frame}



%%%%%%%%%%%%%%%%%%%%%%%%%%%%%%%%%%%%%%%%%%%%%%%%%%%%%%%%%%%%%%%%%%%
\begin{frame}{SVM: Ableiten nach einem Vektor}
\footnotesize
\begin{columns}
\column{0.60\textwidth}
\par\bigskip
Wir suchen das Minimum von $L$ $\Rightarrow$ Ableitung $=0$ setzen.\\
\medskip
Da $L$ mehrdimensional ist, müssen wir partiell ableiten: nach $w_0$ und $\textbf{w}$.\\
\medskip
Die Ableitung nach $\mathbf{w}$ ist der \textbf{Gradient}: man leitet nach
jeder Komponente $w_1, \dots, w_n$ ab und fasst sie zum Vektor zusammen.\\
\medskip
Zwei Regeln genügen für die SVM:
\begin{itemize}
  \item $\frac{1}{2}\lVert\mathbf{w}\rVert^2 = \frac{1}{2}\sum_k w_k^2$
        liefert komponentenweise $w_k$:
        \[
          \frac{\partial}{\partial \mathbf{w}}\Bigl(\tfrac{1}{2}\lVert\mathbf{w}\rVert^2\Bigr)
          = \mathbf{w}
        \]
  \item Das Skalarprodukt $\mathbf{w}\cdot\mathbf{x} = \sum_k w_k x_k$ ist
        linear und liefert:
        \[
          \frac{\partial}{\partial \mathbf{w}}(\mathbf{w}\cdot\mathbf{x})
          = \mathbf{x}
        \]
\end{itemize}

\column{0.35\textwidth}

\begin{block}{\footnotesize Analogie zum Skalaren}
\[
  \tfrac{d}{dx}\bigl(\tfrac{1}{2}x^2\bigr) = x
\]
\[
  \tfrac{d}{dx}(a\,x) = a
\]

\vspace{0.3em}
Der Faktor $\tfrac{1}{2}$ kürzt beim Ableiten die $2$ aus dem Exponenten
weg.
\end{block}
\end{columns}
\source{\cite{scholkopf2002learning}}

\note{
$ \frac{\partial}{\partial \mathbf{w}}\Bigl(\tfrac{1}{2}\lVert\mathbf{w}\rVert^2\Bigr)$ ist partielle Ableitung des Terms in der Klammer.
}

\end{frame}

\begin{frame}{SVM: Ableitungen und optimales $\mathbf{w}$}
\small
\begin{columns}
\column{0.60\textwidth}
Wir leiten die Lagrange-Funktion ab. Konstante Terme und $w_0$ hängen nicht
von $\mathbf{w}$ ab und fallen weg:
\[
  \frac{\partial L}{\partial \mathbf{w}} =
  \mathbf{w} - \sum_{i=1}^{n} \lambda_i y_i \mathbf{x}_i
\]
\[
  \frac{\partial L}{\partial w_0} =
  -\sum_{i=1}^{n} \lambda_i y_i
\]

\vspace{0.3em}
Im \textbf{Optimum} müssen beide Ableitungen \textbf{null} sein
(Stationarität). Setzt man sie null, ergeben sich:
\[
  \mathbf{w}^{*} = \sum_{i=1}^{n} \lambda_i y_i \mathbf{x}_i
  \qquad
  \sum_{i=1}^{n} \lambda_i y_i = 0
\]

\column{0.35\textwidth}
\begin{block}{\small Warum null?} 
    An einem Optimum ist die Steigung waagerecht -- die
Ableitung wird null. Das ist die \textbf{Stationaritätsbedingung} der KKT.
\end{block}

\vspace{0.3em}
\begin{blockC2}{\small Kernaussage} 
 $\mathbf{w}^{*}$ ist eine \textbf{gewichtete Summe}
der Trainingsvektoren; nur Support Vectors tragen bei ($\lambda_i \neq 0$).
\end{blockC2}
\end{columns}
\source{\cite{scholkopf2002learning}}

\note{
Als Folge aus $\sum_{i=1}^{n} \lambda_i y_i = 0$ heben sich die mit
$\lambda_i$ gewichteten Klassenlabels zu null auf. Da $y_i$ nur $+1$
(Klasse A) und $-1$ (Klasse B) annimmt, gilt:
\[
  \sum_{i:\, y_i = +1} \lambda_i = \sum_{i:\, y_i = -1} \lambda_i
\]
Die Gewichte der Support Vectors sind also zwischen beiden Klassen
\textbf{ausbalanciert}. Keine Klasse dominiert die Lage der Trennebene --
beide Seiten ziehen mit gleichem Gesamtgewicht.
}

\end{frame}

%%%%%%%%%%%%%%%%%%%%%%%%%%%%%%%
\begin{frame}{SVM: Herleitung der dualen Formulierung (1)}
\begin{columns}

\column{0.60\textwidth}
\footnotesize
\begin{itemize}
  \item Ausgangspunkt ist unsere Lagrange-Funktion:
        \[
          L_d = \tfrac{1}{2}\mathbf{w}\cdot\mathbf{w}
          - \sum_{i=1}^{n}\lambda_i \cdot y_i \cdot (w_0 + \mathbf{w}\cdot\mathbf{x}_i)
          + \sum_{i=1}^{n}\lambda_i
        \]
  \item Wir spalten den mittleren Term auf:
        \[
          L_d = \tfrac{1}{2}\mathbf{w}\cdot\mathbf{w}
          - w_0 \underbrace{\sum_{i=1}^{n}\lambda_i \cdot y_i}_{=\,0}
          - \mathbf{w}\cdot\sum_{i=1}^{n}\lambda_i \cdot y_i \mathbf{x}_i
          + \sum_{i=1}^{n}\lambda_i
        \]
  \item Der Term mit $w_0$ \textbf{fällt weg}, da
        $\sum_{i=1}^{n}\lambda_i \cdot y_i = 0$ gilt.
\end{itemize}

\column{0.35\textwidth}
\footnotesize
\textbf{Was passiert hier?} Das Skalarprodukt im mittleren Term wird
ausmultipliziert, sodass $w_0$ und $\mathbf{w}$ getrennt vorkommen. So wird
sichtbar, dass der $w_0$-Term dank der Stationaritätsbedingung
verschwindet.
\end{columns}
\source{\cite{scholkopf2002learning}}
\end{frame}

\begin{frame}{SVM: Herleitung der dualen Formulierung (2)}
\begin{columns}[b]

\column{0.60\textwidth}
\footnotesize
\begin{itemize}
  \item Es bleibt:
        \[
          L_d = \tfrac{1}{2}\mathbf{w}\cdot\mathbf{w}
          - \mathbf{w}\cdot\sum_{i=1}^{n}\lambda_i \cdot y_i \mathbf{x}_i
          + \sum_{i=1}^{n}\lambda_i
        \]
  \item Einsetzen von
        $\mathbf{w}^{*} = \sum_{i=1}^{n}\lambda_i \cdot y_i \mathbf{x}_i$
        in \textbf{beide} $\mathbf{w}$-Terme:
        \[
          L_d = \tfrac{1}{2}
          \Bigl(\sum_{i}\lambda_i \cdot y_i \mathbf{x}_i\Bigr)\!\cdot\!
          \Bigl(\sum_{j}\lambda_j \cdot y_j \mathbf{x}_j\Bigr)
          - \Bigl(\sum_{i}\lambda_i \cdot y_i \mathbf{x}_i\Bigr)\!\cdot\!
          \Bigl(\sum_{j}\lambda_j \cdot y_j \mathbf{x}_j\Bigr)
          + \sum_{i}\lambda_i
        \]
  \item Der erste Term ist die \textbf{Hälfte} des zweiten; zusammengefasst
        bleibt \textbf{ein negativer} Doppelsummen-Term:
        \[
          L_d = -\tfrac{1}{2}\sum_{i=1}^{n}\sum_{j=1}^{n}
          \lambda_i \lambda_j \cdot y_i \cdot y_j\,(\mathbf{x}_i\!\cdot\mathbf{x}_j)
          + \sum_{i=1}^{n}\lambda_i
        \]
\end{itemize}

\column{0.35\textwidth}
\begin{blockC1}{\footnotesize Duales Problem}
\footnotesize
$L_d$ hängt nur noch von $\boldsymbol{\lambda}$ ab und wird bezüglich
$\boldsymbol{\lambda}$ \textbf{maximiert}. Die Daten treten nur als
\textbf{Skalarprodukte} $\mathbf{x}_i\!\cdot\mathbf{x}_j$ auf.
\end{blockC1}
\end{columns}
\source{\cite{scholkopf2002learning}}
\end{frame}


\begin{frame}{SVM: Das duale Problem}
\begin{columns}

\column{0.47\textwidth}
\footnotesize
\begin{itemize}
  \item Das duale Problem der SVM lautet:
        \[
          \max_{\boldsymbol{\lambda}} \; -\tfrac{1}{2}\sum_{i=1}^{n}\sum_{j=1}^{n}
          \lambda_i \lambda_j \cdot y_i \cdot y_j\,(\mathbf{x}_i\!\cdot\mathbf{x}_j)
          + \sum_{i=1}^{n}\lambda_i
        \]
  \item unter den Nebenbedingungen
        $\lambda_i \geq 0$ und
        $\sum_{i=1}^{n}\lambda_i \cdot y_i = 0$.
  \item Es ist ein \textbf{quadratisches Problem} in den $\lambda_i$.
  \item Gelöst wird es \textbf{numerisch}, meist mit \textbf{SMO}
        (Sequential Minimal Optimization).
  \item Aus $\boldsymbol{\lambda}^{*}$ berechnet man $\mathbf{w}^{*}$ und
        $w_0^{*}$.
\end{itemize}

\column{0.48\textwidth}
\begin{blockDef}{\footnotesize Duales Problem}
\footnotesize
\begin{itemize}
  \item Ein \textbf{duales Problem} ist eine alternative Formulierung des
        ursprünglichen (\textbf{primalen}) Problems.
  \item \textbf{Primal:} Optimierung über $\mathbf{w}$ und $w_0$ unter den
        Trennungs-Nebenbedingungen.
  \item \textbf{Dual:} Über die Lagrange-Funktion werden $\mathbf{w}$ und
        $w_0$ eliminiert; optimiert wird nur noch über
        $\boldsymbol{\lambda} = (\lambda_1, \dots, \lambda_n)$.
  \item Beide führen zur \textbf{selben Lösung} (starke Dualität bei
        konvexen Problemen).
\end{itemize}
\end{blockDef}
\end{columns}
\source{\cite{scholkopf2002learning}}
\note{Die konkrete Lösung ist uns egal. Wir brauchen für später nur die Duale formulieren. Wegen dem Skalarprodukt}
\end{frame}

%%%%%%%%%%%%%%%%%%%%%%%%%%%%%%%%


\begin{frame}{SVM: Effizienz durch Support Vectors}
\begin{columns}

\column{0.47\textwidth}
\centering
\includesvg[width=\linewidth]{figures/svm_synthetisch_supportvektoren.svg}

\column{0.48\textwidth}
\footnotesize
\begin{blockC1}{Effizienz}
\begin{itemize}
  \item Nur für \textbf{Support Vectors} ist $\lambda_i \neq 0$.
  \item Nur ein \textbf{Bruchteil} der Daten geht in die Rechnung ein.
  \item Der Aufwand hängt an der \textbf{Anzahl der Support Vectors}, nicht
        an der Gesamtmenge.
\end{itemize}
\end{blockC1}

\vspace{0.3em}
{\footnotesize\itshape Nicht linear trennbare Probleme lösen Soft-Margin
und Kernel-Ansätze.}
\end{columns}
\source{\cite{scholkopf2002learning}, SVM\_Linear\_Synthetisch.ipynb}

\end{frame}




\makequizpage{\QUIZURL}{\QUIZQRCODE}{\QUIZIMAGE}
\makequestionspage{\FRAGENIMAGE}